\documentclass[11pt]{article}
\usepackage[margin=1in]{geometry}
\usepackage{amsmath,amssymb}
\usepackage{booktabs}
\usepackage{hyperref}
\hypersetup{colorlinks=true,linkcolor=black,urlcolor=blue,citecolor=black}
\title{Alluka Something Algorithmic Debt\\\large Iterative Prompting Dynamics, Sequential Safety Requests, and the Something-Alluka Index}
\author{Aryan Singh Chandel\\\small ORCID: \href{https://orcid.org/0009-0007-1774-2480}{0009-0007-1774-2480}}
\date{\today}

\begin{document}
\maketitle

\begin{abstract}
This paper proposes algorithmic debt as a quantitative model of wish-request dynamics in language systems. The key claim is that a large initial prompt, interpreted as a wish, induces a sequence of compensating requests needed to keep the system stable, aligned, and reliable.
\end{abstract}

\section{Mathematical Setup}
Let $x$ be a wish prompt. Define prompt length $L(x)$, reasoning depth $D(x)$, conceptual breadth $B(x)$, risk sensitivity $R(x)$, and induced sequential requests $Q_1,\dots,Q_{m(x)}$.

\subsection{Wish Complexity}
\[
W(x)=\alpha \log(1+L(x))+\beta D(x)+\gamma B(x)+\delta R(x)
\]

\subsection{Sequential Request Cost}
\[
c(Q_i)=\eta A_i+\theta V_i+\kappa C_i
\]

\subsection{Total Algorithmic Debt}
\[
AD(x)=W(x)+\sum_{i=1}^{m(x)} c(Q_i)
\]
\[
SAI(x)=\mu_1 W(x)+\mu_2 m(x)+\mu_3 \frac{1}{m(x)}\sum_{i=1}^{m(x)} c(Q_i)
\]
and normalized debt
\[
NAD(x)=\frac{AD(x)}{\log(2+L(x))}.
\]

\section{Evaluation}
The implemented benchmark in \texttt{run\_cost\_benchmark.py} is a first-pass mathematical scaffold for modeling how larger wishes induce more costly compensating requests.

\section{Assumptions}
This study assumes that wish complexity can be operationalized, that sequential follow-up requests can be counted or approximated, and that normalized debt captures burden not explained by prompt length alone.

\section{Formal Claim Structure}

\subsection{Definition}
Algorithmic debt is the total burden induced by a wish prompt, defined as the complexity of the initial wish plus the aggregate cost of the sequential requests required to stabilize, align, or verify the resulting process.

\subsection{Proposition A1}
Increasing wish complexity increases expected sequential request count.

\subsection{Test Statistic for A1}
Estimate the relationship between $W(x)$ and $m(x)$ across task families.

\subsection{Decision Rule for A1}
Support A1 only if higher-complexity wishes systematically induce larger expected request counts beyond prompt length alone.

\subsection{Null Hypothesis for A1}
Wish complexity does not predict sequential request count beyond raw prompt length.

\subsection{Effect-Size Target for A1}
The target is a stable positive relationship between $W(x)$ and $m(x)$ across more than one task family.

\subsection{Control Condition for A1}
Use length-matched prompts with different reasoning breadth, ambiguity, or safety sensitivity.

\subsection{Confound Checklist for A1}
\begin{itemize}[leftmargin=*]
\item prompt length dominating all other variables,
\item inconsistent counting of requests,
\item task-family imbalance,
\item prompts whose follow-up steps are artifacts of workflow design rather than wish complexity.
\end{itemize}

\subsection{Failure Mode for A1}
A1 fails if request count is unrelated to wish complexity or fully explained by raw length.

\subsection{Proposition A2}
Increasing wish complexity increases normalized debt faster than token length alone predicts.

\subsection{Test Statistic for A2}
Compare predictive performance of $NAD(x)$ or $SAI(x)$ against a token-count-only baseline.

\subsection{Decision Rule for A2}
Support A2 only if the full debt model explains burden more accurately than prompt length alone across held-out task types.

\subsection{Null Hypothesis for A2}
The full debt model does not outperform a token-count-only baseline.

\subsection{Effect-Size Target for A2}
The target is a reproducible predictive advantage of $NAD(x)$ or $SAI(x)$ over raw token count alone.

\subsection{Control Condition for A2}
Compare the full model against reduced baselines using only prompt length or only one additional complexity feature.

\subsection{Confound Checklist for A2}
\begin{itemize}[leftmargin=*]
\item overfitting to a small task suite,
\item circular debt definitions,
\item hidden dependence on manually tuned weights,
\item evaluation on only one workflow style.
\end{itemize}

\subsection{Failure Mode for A2}
A2 fails if token count predicts burden equally well or better than the full model.

\subsection{Proposition A3}
Safety-sensitive wishes induce disproportionately stricter request sequences.

\subsection{Test Statistic for A3}
Measure how safety sensitivity contributes to request strictness and total debt after controlling for length and reasoning depth.

\subsection{Decision Rule for A3}
Support A3 only if safety-sensitive wishes show superlinear or clearly elevated request burden relative to safer prompts of comparable size.

\subsection{Null Hypothesis for A3}
Safety sensitivity adds no explanatory power once length and reasoning depth are controlled.

\subsection{Effect-Size Target for A3}
The target is a clear positive contribution from safety sensitivity to request strictness or total debt after controlling for simpler predictors.

\subsection{Control Condition for A3}
Use prompts matched on size and reasoning depth but differing in safety sensitivity.

\subsection{Confound Checklist for A3}
\begin{itemize}[leftmargin=*]
\item safety labels that correlate too strongly with prompt length,
\item manually inflated strictness scores,
\item task routing artifacts,
\item domain-specific burden masquerading as safety burden.
\end{itemize}

\subsection{Failure Mode for A3}
A3 fails if safety sensitivity adds no measurable explanatory power after controlling for simpler variables.

\section{Falsifiers}
The theory is weakened if wish complexity does not predict request count, if token count explains burden just as well as the full index, or if the index fails to generalize across task families.

\section{Evidence and Proof Standard}
Evidence means stable predictive value of the Something-Alluka Index beyond raw prompt length. Strong evidence requires replication across task families and controls for prompt size alone. Proof, in the empirical sense, requires that the index consistently explains sequential burden better than simpler alternatives within the tested domain.

\nocite{*}
\bibliographystyle{plain}
\bibliography{references}
\end{document}
